\section{The whole-epitope antibody set, built at native 103}\label{sec:c1set}
Having settled on RFdiffusion3 (Section~\ref{sec:rfd}), we built the antibody design set it feeds. For
each of the \num{56} known-antibody mAb epitopes (the three non-standard cases of
Section~\ref{sec:dp4c6} excluded) we scaffold the \emph{whole} epitope --- all of its islands held
together in their native geometry, the way the antibody sees it --- into a constant-length construct.
This is the C1 component of the shipped library (Section~\ref{sec:dp4c1}) and the pool from which the
metric-space sample (C5) and the mutant controls (C6) are drawn, so its properties set the ceiling for
three of the four antibody components.

\paragraph{Native 103, not a truncated 104.} The construct length is fixed at \num{103} residues, the
PepSeq maximum. Lawson's original whole-epitope contigs were \num{104} (Section~\ref{sec:rfd}), and our
first pass reproduced them and then trimmed the resulting \num{104}-mers to \num{103} at assembly. We
discarded that route in favor of regenerating the set natively at \num{103}. The DP3 assay we build on
itself included DP2 \num{104}-mers truncated to \num{103}, and these showed generally weaker binding
signal --- most plausibly assay run-to-run variation, but enough to make a post-hoc single-residue
truncation a liability we prefer not to carry on the components that ride on this pool. Instead we take
Lawson's exact whole-epitope contigs and remove one \emph{scaffold} residue
(\texttt{build\_whole\_epitope\_designs.py}) --- the larger terminal flank, or the largest inter-island
spacer when the islands sit flush against both termini --- preserving every epitope residue and his
inter-island-spacing sweep, so the set differs from the \num{104} version only in length. That the
length change is metric-neutral is shown directly (Section~\ref{sec:rfd}, Fig.~\ref{fig:len}). Building
at \num{103} natively also dissolves the one case the trim could not handle: the epitope whose two
islands are flush against both termini (\texttt{3ux9\_1P}) simply has its inter-island spacer shortened,
keeping both islands intact, with no design dropped.

\paragraph{The 103-mer run.} This gives \num{2206} contigs over the \num{56} epitopes, each sampled at \num{8}
RFdiffusion3 backbones $\times$ \num{8} ProteinMPNN sequences $=\num{64}$ designs and pushed through
AlphaFold3, of which \num{140716} returned a prediction. We extract every per-design metric once into a
single table --- the C1 analog of \texttt{dp2} (\texttt{scripts/stage05\_extract\_metrics.py}) --- and
make all downstream choices from it. The four filters pass \num{727}/\num{140716} designs ($0.52\%$),
in line with the \num{104}-mer set on the same epitopes ($0.35\%$) and with Lawson's RFD1 ($0.56\%$):
the sub-percent yield expected for whole-epitope scaffolding (Section~\ref{sec:intro}), and unchanged by
the move to \num{103}.

\paragraph{Splitting by island count.} Splitting the set by island count (Fig.~\ref{fig:c1island}) recovers the
expected asymmetry. The \num{13} single-island epitopes are cleaner on every metric (median epitope RMSD
\num{2.40} vs \num{4.99}~\si{\angstrom}, median PAE \num{9.8} vs \num{12.8}) and pass at $2.79\%$ against
$0.50\%$ for the \num{43} two-island epitopes: holding two separated patches in their native arrangement
is simply harder than holding one (Section~\ref{sec:islands}). This is a cleaner result than the DP3
set, where the binary pass rate flipped on two outlier epitopes; here the single-island advantage is
consistent between the metric distributions and the pass rate.

\begin{figure}[h]\centering
\figorbox{metrics_103_island_split.png}{1.0}
\caption{Native-103 whole-epitope set: four-filter metric distributions split by epitope island count
(1 island, blue; 2 islands, red), density-normalized over designs with an AlphaFold3 result.
Single-island designs sit at better values on every metric and pass more often ($2.79\%$ vs $0.50\%$).
Dotted lines are group medians; dashed lines the filter thresholds. Regenerated by
\texttt{episcaf\_analysis/viz/plot\_whole\_epitope\_103.py}.}
\label{fig:c1island}
\end{figure}

\paragraph{Selecting the C1 set.} We rank on the re-weighted composite with no hard gate
(Section~\ref{sec:selection}) and keep the top \num{20} per epitope, \num{1120} designs, as C1; because
the antibody is known, the accessibility term is the real AlphaFold3 clash rather than the cylinder
surrogate. Case-encoding these designs, and the C5/C6 rebuilds off this pool, follow in
Section~\ref{sec:dp4}.
